% This is the ADASS_template.tex LaTeX file, 24 Oct 2025.
% It is based on the ASP general author template file, but modified to reflect the specific
% requirements of the ADASS proceedings.

% This is the aspauthor.tex LaTeX file
% Copyright 2021, Astronomical Society of the Pacific Conference Series
% Revision:  29 November 2021

% To compile, at the command line positioned at this folder, type:
% pdflatex ADASS_template
% pdflatex ADASS_template
% This will create a file called ADASS_template.pdf

\documentclass[11pt,twoside]{article}

% Do NOT use ANY packages other than asp2023.
\usepackage{asp2023}

\aspSuppressVolSlug
\resetcounters

% References must all use BibTeX entries in a .bibfile.
% References must be cited in the text using \citet{} or \citep{}.
% Do not use \cite{}.
% See ManuscriptInstructions.pdf for more details
\bibliographystyle{asp2021}

\markboth{Stoehr, F. et al.}{ALMA Spectral Line Search}

\begin{document}

\title{The ALMA spectral line search is now online}

\author{Felix~Stoehr$^1$, Andrea Farago$^1$, Stefan Curiban$^1$, Alisdair Manning$^1$, Jorge~Garcia$^2$, Pei-Ying Hsieh$^3$, Andy Lipnicky$^4$, and Adele~Plunkett$^4$}
\affil{$^1$European Southern Observatory (ESO), Garching, Germany; \email{fstoehr[at]eso[dot]org}}
\affil{$^2$Joint ALMA Observatory (JAO), Chile}
\affil{$^3$National Astronomical Observatory of Japan (NAOJ), Japan}
\affil{$^4$National Radio Astronomy Observatory (NRAO), Charlottesville, USA}

\paperauthor{Felix~Stoehr}{fstoehr (at) eso (dot) org}{}{ESO}{ALMA}{Garching}{}{85748}{Germany}
\paperauthor{Andrea~Farago}{}{}{ESO}{ALMA}{Garching}{}{85748}{Germany}
\paperauthor{Stefan~Curiban}{}{}{ESO}{ALMA}{Garching}{}{85748}{Germany}
\paperauthor{Alisdair~Manning}{}{}{ESO}{ALMA}{Garching}{}{85748}{Germany}
\paperauthor{Jorge~Garcia}{}{}{JAO}{ALMA}{Santiago}{}{19001}{Chile}
\paperauthor{Pei-Ying~Hsieh}{}{0000-0003-2224-4922}{NAOJ}{}{Mitaka}{Tokyo}{181-8588}{Japan}
\paperauthor{Andy~Lipnicky}{}{0000-0002-6667-7773}{NRAO}{}{Charlottesville}{VA}{22903}{USA}
\paperauthor{Adele~Plunkett}{}{0000-0002-9912-5705}{NRAO}{}{Charlottesville}{VA}{22903}{USA}

%\aindex{Stoehr,~F.}
%\aindex{Farago,~A.}
%\aindex{Curiban,~S.}
%\aindex{Manning,~A.}
%\aindex{Garcia,~J.}
%\aindex{Hsieh,~P.-Y.}
%\aindex{Lipnicky,~A.}
%\aindex{Plunkett,~A.}

\begin{abstract}
The ALMA Science Archive (ASA \url{https://almascience.org/aq}) Query interface now allows users to search for lines that have been actually {\it detected} in public ALMA data cubes. To that end, we have run the post-processing softwae ALMA Data Mining Toolkit (ADMIT) \citep{admit} on all ALMA data-cubes, once at native resolution and once binned by a factor 16, and have ingested the metadata into the ASA database. Users can now not only search for molecules and transitions, but also for redshifts and source velocities or line intensities. We also offer previews for the detected lines as created by ADMIT.

While we expect this new functionality to be extremely a very useful addition to the ASA, we do note that it comes with a large number of caveats. This is due to the often uncertain or incorrect source redshifts or velocities, as well as partia incompleteness of the information presented in the case that, for example, not all spectral windows in the raw data also were imaged by ALMA. We ask users to proceed with utmost care and use their own scientific judgement.
\end{abstract}

%\ssindex{archives!individual!ALMA Science Archive}
%\ooindex{ADMIT, ascl:XXXXX}

\section{Introduction}
The amount of astronomical data is growing exponentially \citep[e.g.][]{stoehradass19} resulting in an exponential increase of data available per astronomer. Not data but astronomers are becoming - or have already become - the rare resource in astronomy \citep{stoehradass15}. For science archives like the ALMA Science Archive, this evolution means that the next frontier is to make the actual content of the data searchable \citep{durandadass17}, in addition to the searches offered on observational metadata. ALMA has taken a few steps into this direction \citep{stoehrmessenger23}: users can search for similar abstracts of publications and projects on the ASA query interface using a state-of-the-art text similarity method. Additionally, we offer preview images that show lines identified through the ALMA Data Mining Toolkit (ADMIT) \citep{admit}. Recently we have added image similarity search (\citet{stoehradass23}, \citet{stoehradass25}).

With the line search described in this work, we present our most recent effort to make the content of ALMA observations searchable by extracting the metadata from the ADMIT runs, ingesting them into the ASA database and augmenting the web search interface accordingly.

\section{Running ADMIT}
\label{admit}
We have run ADMIT on all FITS files in the ASA of which resulted in about 695,000 ADMIT files of ALMA data cubes. We ran two ADMIT proceesings per cube: once at the native resolution of the data-cube and once where the cube is binned by a factor of 16 in the spectral direction. Running ADMIT on a binned cube increases the signal to noise ratio and allows the detection of broad weak lines, relevant for extragalactic observations.

In addition to the FITS cube, ADMIT also takes a redshift or the vlsr, the velocity of the source in the local standard of rest. We obtain the initial values of the redshift or vlsr from the information the Principal Investigator (PI) provides in the ALMA Observing Tool (OT).

Unfortunately, these values are only of mediocre quality for a large number of reasons. In some cases, PIs set the velocity of their source to zero and adjust the spectral windows to observe even if the source does have a vlsr quite different from zero. In other cases, the user provides the vlsr of the not of the gravitational lensed object but that of the lens. Or the PI adjusts the spectral window of an observation of a field of galaxies so that all relevant lines of the galaxies fall into the spectral window, but the vlsr does not correspond to the vlsr of any of the objects. In a very large number of cases no redshift or vlsr are provided in the first place. In some cases, the redshift of the source is not known and the goal of the ALMA observations is to determine that redshift.

To at least somewhat improve on the situation, we apply the following algorithm to the redshift or vlsr values we are given by the PIs in the OT: 1) If a velocity was given by the PI, in the OT, that velocity is used. 2) If no velocity was given, we look for other existing ALMA observations of the same source and use the velocity if a value is available. 3) If not, we try to match the observed source with a source from SIMBAD and NED and use the velocity from those services. 4) If none is available, we check to see if the PI has provided lines they were looking for in the OT and compute the velocities from that information. This process provides guesses for velocities for about 80\% of all FITS files. In a significant fraction of the cases, however, the values do not correspond to what a scientist would be able to determine upon careful investigation.

We ingest the ADMIT products into the ALMA Science Archive so that they can be downloaded by PIs and archival researchers.

In a separate process, we then obtain all metadata files from the Archive, extract the lines found by ADMIT including their parameters and ingest them into the ALMA Science Archive database. At the same time, we also extract the preview images of the line emission, change the color table to the one which is used by the AQ (matplotlib's viridis colour table) and ingest those files alongside the metadata. The metadata and images then get served through the ALMA Science Archvie Query interface (AQ hereafter).

\section{Search Interface}
\articlefigure[width=0.9\textwidth]{P806_01.eps}{figure_01}{ASA query interface showing the new input fields related to metadata of detected spectral lines, as well as at the bottom the line view of the result table and a spectral line preview.}

The search interface now offers the new search input fields Rest Frequency, Line, Redshift, VLSR, SNR, FWHM and Peak Intensity. When typing into the Line field a molecule or transition, the interface displays a drop-down menu of the subset of the 92,000 lines of the official ALMA-subset of splatalogue\footnote{\url{https://splatalogue.online}} that match the user-entry - sorted by the popularity of the lines in the ALMA observations of the last 15 years. This very effectively allows users to select the lines they are interested in.

In addition to new search parameters, the interface also provides a new tab in the result table. Allowing to display the results of the search not only viewing by observations, projects or publications, but now also in a way where each displayed row corresponds to one detected line. The interface there provides the source name, molecule, transition, FWHM of the line, redshift, VLSR, rest frequency, peak intensity, SNR and the name of the FITS file in which the line was detected. Icons allow users to jump to the place on the sky of the observation, to see the spectral coverage in the frequency viewer, find similar projects based on the proposal text, download the FITS file, find objects that are morphologically similar and - last not least - see the preview of the line emission.

\section{Caveats}
While the AQ provides only highly curated and correct values and estimates for users to search on, the values related to the line search, come with a number of caveats.

\begin{enumerate}
\item As mentioned in Section~\ref{admit} the quality of the redshift or vlsr values available is only mediocre which often results in assigning wrong labels to identified lines. All values and labels of questionable accuracy are therefore presented in a different colour - orange in this case - to make a clear distinction from the values that can be trusted.
\item In addition - and in contrast to the results presented in the observations, project and publications tabs, the information in the lines tab is also highly incomplete. ALMA does not produce products for all spectral windows and spectral channels in the first place. Also, the automatic extraction of lines is limited to a SNR of 6, and only the lines for the most prominent molecules are identified.
\end{enumerate}

To ensure that really each user unerstands and acknowledges the caveats, a pop-up window is presented as soon as line data is searched for or the result table is displayed. Moreover, we have removed the programmatic access through download of VO Tables or through TAP for those metadata so that users or AI agents are prevented from accidentially accessing the line metadata  without having understood the limitations.

\acknowledgements
ALMA is a partnership of ESO (representing its member sta-tes), NSF (USA) and NINS (Japan), together with NRC (Canada), NSTC and ASIAA (Taiwan), and KASI (Republic of Korea), in cooperation with the Republic of Chile. The Joint ALMA Observatory is operated by ESO, AUI/NRAO and NAOJ.



\bibliography{P806}

% if we have space left, we might add a conference photograph here. Leave commented for now.
% \bookpartphoto[width=1.0\textwidth]{foobar.eps}{FooBar Photo (Photo: Any Photographer)}

\end{document}
